% Options for packages loaded elsewhere
\PassOptionsToPackage{unicode}{hyperref}
\PassOptionsToPackage{hyphens}{url}
\PassOptionsToPackage{dvipsnames,svgnames,x11names}{xcolor}
\PassOptionsToPackage{space}{xeCJK}
\documentclass[
]{article}
\usepackage{xcolor}
\usepackage[margin=2.5cm]{geometry}
\usepackage{amsmath,amssymb}
\setcounter{secnumdepth}{-\maxdimen} % remove section numbering
\usepackage{iftex}
\ifPDFTeX
  \usepackage[T1]{fontenc}
  \usepackage[utf8]{inputenc}
  \usepackage{textcomp} % provide euro and other symbols
\else % if luatex or xetex
  \usepackage{unicode-math} % this also loads fontspec
  \defaultfontfeatures{Scale=MatchLowercase}
  \defaultfontfeatures[\rmfamily]{Ligatures=TeX,Scale=1}
\fi
\usepackage{lmodern}
\ifPDFTeX\else
  % xetex/luatex font selection
    \setmainfont[]{DejaVu Sans}
    \setmonofont[]{DejaVu Sans Mono}
  \ifXeTeX
    \usepackage{xeCJK}
    \setCJKmainfont[]{Noto Sans CJK SC}
          \fi
  \ifLuaTeX
    \usepackage[]{luatexja-fontspec}
    \setmainjfont[]{Noto Sans CJK SC}
  \fi
\fi
% Use upquote if available, for straight quotes in verbatim environments
\IfFileExists{upquote.sty}{\usepackage{upquote}}{}
\IfFileExists{microtype.sty}{% use microtype if available
  \usepackage[]{microtype}
  \UseMicrotypeSet[protrusion]{basicmath} % disable protrusion for tt fonts
}{}
\makeatletter
\@ifundefined{KOMAClassName}{% if non-KOMA class
  \IfFileExists{parskip.sty}{%
    \usepackage{parskip}
  }{% else
    \setlength{\parindent}{0pt}
    \setlength{\parskip}{6pt plus 2pt minus 1pt}}
}{% if KOMA class
  \KOMAoptions{parskip=half}}
\makeatother
\usepackage{longtable,booktabs,array}
\usepackage{calc} % for calculating minipage widths
% Correct order of tables after \paragraph or \subparagraph
\usepackage{etoolbox}
\makeatletter
\patchcmd\longtable{\par}{\if@noskipsec\mbox{}\fi\par}{}{}
\makeatother
% Allow footnotes in longtable head/foot
\IfFileExists{footnotehyper.sty}{\usepackage{footnotehyper}}{\usepackage{footnote}}
\makesavenoteenv{longtable}
\setlength{\emergencystretch}{3em} % prevent overfull lines
\providecommand{\tightlist}{%
  \setlength{\itemsep}{0pt}\setlength{\parskip}{0pt}}
\usepackage{bookmark}
\IfFileExists{xurl.sty}{\usepackage{xurl}}{} % add URL line breaks if available
\urlstyle{same}
\hypersetup{
  colorlinks=true,
  linkcolor={Maroon},
  filecolor={Maroon},
  citecolor={Blue},
  urlcolor={Blue},
  pdfcreator={LaTeX via pandoc}}

\author{}
\date{}

\begin{document}

\section{Three-Channel Equivalence in the Unified Framework of
Representation, Logic, and Intuition --- Construction Earns
Intuition}\label{three-channel-equivalence-in-the-unified-framework-of-representation-logic-and-intuition-construction-earns-intuition}

\section{统一框架 · 三通道等价 ·
构造获得直觉}\label{ux7edfux4e00ux6846ux67b6-ux4e09ux901aux9053ux7b49ux4ef7-ux6784ux9020ux83b7ux5f97ux76f4ux89c9}

\begin{quote}
\textbf{Draft v0.1 (English)} \textbar{} 2026-08-11 Target: TMLR / COLM
/ Cognition (theory + empirical) \textbf{Framework position}: Part of
the \emph{Unified Framework of Representation, Logic, and Intuition}
(manifesto:
\texttt{Unified\_Framework\_Representation\_Logic\_Intuition.md}). This
paper establishes the \textbf{construction→intuition} cognition:
construction, intuition, and formal rewriting are three equivalent
execution channels of one system; intuition is compiled construction.
Companion: P1 (form→construction, representation + generalization
theory), P2 (construction→intuition method, compilation). Evidence:
exp50 (three-channel consistency), exp41 (symbol permutation), exp80
(step-skipping), exp53 (convergence). Single official run\_exp judge.
\end{quote}

\begin{center}\rule{0.5\linewidth}{0.5pt}\end{center}

\textbf{Repository:
\href{https://github.com/YuchenWang-ai/Three-Channel-Equivalence-Construction-Intuition-and-Formal-Rewriting-in-One-Neural-System}{YuchenWang-ai/Three-Channel-Equivalence-Construction-Intuition-and-Formal-Rewriting-in-One-Neural-System}}

\subsection{Abstract}\label{abstract}

The century-old split in the foundations of mathematics --- intuitionism
(a mathematical object exists iff it can be constructed), structuralism
(an object is its position in a structure), and formalism (mathematics
is a sign game) --- is mirrored in a single neural system as three
execution channels for the same mathematical object. We show that, in a
token-native representation, the \emph{construction} channel
(definition-driven unfolding), the \emph{intuition} channel (compiled
weight execution), and the \emph{formal} channel (definition-driven
rewriting) are pairwise semantically equivalent: on a 2000-digit carry
chain all three produce identical results (1.000). We argue this is the
neural implementation of Computational Trinitarianism (Curry--Howard:
programs = proofs = types), with the crucial difference that the
equivalence holds between \emph{execution channels of one system} rather
than between logics. We develop the epistemology: intuition is the
reasoning fast path --- hard to obtain (config-dependent epochs), but
easy to generalize by skipping construction steps; symbol permutation
invariance (0.987) and zero-decay step-skipping (2000 digits, radix 60,
100\% anonymization) establish that the intuition channel carries
structure, not symbol semantics. Stance: anti-Platonist --- intuition is
compiled construction, not innate. Within the Unified Framework, this
paper is the \textbf{construction→intuition} cognition leg: correct
construction yields intuition (compiled construction), and correct
generalization compiles into flashing intuition (the fast reasoning
path).

\subsection{1 Introduction}\label{introduction}

\textbf{The split.} In the foundations of mathematics, three answers
compete for the question ``what is a mathematical object?'' Brouwer's
intuitionism: an object exists iff it can be \emph{constructed}.
Bourbaki's structuralism: an object is its \emph{position in a
structure}. Hilbert's formalism: mathematics is a \emph{sign game}; an
object is its role in a rewriting system. The dispute is usually
presented as philosophical; in this paper we observe that it is also
\emph{architectural} --- the three positions correspond to three ways a
computing system can realize the same object.

\textbf{The unification.} A single neural system over a token-native
representation (P1) exhibits all three channels: - \emph{Construction}
(intuitionism): unfold the definition chain, step by step --- existence
by construction; - \emph{Intuition} (structuralism): compiled weights
position the object directly --- existence by location in the learned
structure; - \emph{Formal} (formalism): the definition-driven rewriting
engine --- existence by the sign-manipulation rules.

We show experimentally that the three channels are pairwise semantically
equivalent on the same inputs (exp50: construction = formal exactly;
intuition 1.000 on a 2000-digit carry chain). This is the neural
implementation of Computational Trinitarianism --- Curry--Howard's
programs = proofs = types --- but with a decisive difference: the
equivalence is \emph{within one system's execution channels}, not
between distinct logical systems.

\textbf{Why this matters.} If construction, intuition, and formal
rewriting agree on the same object, then the philosophical positions are
not alternatives to be argued about but \emph{facets of a single
realization}. The disagreement over ``what an object is'' becomes a
disagreement over \emph{which channel to look through} --- and all
channels answer the same way. The engineering content is the compilation
relation: construction can be compiled into intuition (P2), and both
agree with rewriting.

\textbf{Contributions.}

\begin{itemize}
\tightlist
\item
  \textbf{C1 (Epistemology) --- Intuition is compiled construction.} The
  intuition channel is not a third mysterious faculty; it is the
  construction process compiled into weights. Anti-Platonist: no innate
  structure, no a priori object.
\item
  \textbf{C2 (Empirics) --- Three-channel equivalence.} Construction =
  formal exactly, and intuition 1.000 on the same inputs, including a
  2000-digit carry chain (exp50).
\item
  \textbf{C3 (Empirics) --- Intuition carries structure, not symbols.}
  Symbol permutation invariance (0.987, exp41) and zero-decay
  step-skipping (exp80) show the intuition channel realizes the
  structure of the objects, not their signifiers.
\item
  \textbf{C4 (Epistemology) --- Acquisition and step-skipping are
  decoupled.} Intuition is hard to obtain (config-dependent epochs; 3
  simple / 8 full, exp53) but easy to generalize by skipping
  construction steps --- the two facts are not in tension.
\end{itemize}

\subsection{2 Three Positions, Three
Channels}\label{three-positions-three-channels}

\begin{longtable}[]{@{}
  >{\raggedright\arraybackslash}p{(\linewidth - 8\tabcolsep) * \real{0.2000}}
  >{\raggedright\arraybackslash}p{(\linewidth - 8\tabcolsep) * \real{0.2000}}
  >{\raggedright\arraybackslash}p{(\linewidth - 8\tabcolsep) * \real{0.2000}}
  >{\raggedright\arraybackslash}p{(\linewidth - 8\tabcolsep) * \real{0.2000}}
  >{\raggedright\arraybackslash}p{(\linewidth - 8\tabcolsep) * \real{0.2000}}@{}}
\toprule\noalign{}
\begin{minipage}[b]{\linewidth}\raggedright
Position
\end{minipage} & \begin{minipage}[b]{\linewidth}\raggedright
Object is
\end{minipage} & \begin{minipage}[b]{\linewidth}\raggedright
Access
\end{minipage} & \begin{minipage}[b]{\linewidth}\raggedright
Our channel
\end{minipage} & \begin{minipage}[b]{\linewidth}\raggedright
Execution
\end{minipage} \\
\midrule\noalign{}
\endhead
\bottomrule\noalign{}
\endlastfoot
Intuitionism (Brouwer) & a construction & construct it & Construction &
unfold definition chain, step by step \\
Structuralism (Bourbaki) & a position & locate it & Intuition & compiled
weights place it directly \\
Formalism (Hilbert) & a sign game & play it & Formal & definition-driven
rewriting \\
\end{longtable}

The three channels share a single definition source: the token-native
representation (P1 §4). Construction and formal read the same
definitions; intuition is a compiled image of the same definitions.
Equivalence is therefore not a coincidence of three independent systems
but a property of one system viewed through three channels.

\textbf{Computational Trinitarianism.} Curry--Howard identifies
programs, proofs, and types as the same thing under three presentations.
Our claim is analogous: construction, intuition, and formal rewriting
are the same computation under three presentations --- \emph{within one
system}. The difference from Curry--Howard is the locus of equivalence:
not ``program = proof = type'' as distinct entities, but ``one object's
three execution paths agree.''

\subsection{3 Intuition as Compiled
Construction}\label{intuition-as-compiled-construction}

\textbf{The claim.} The intuition channel is the construction process
compiled into weights. Training (P1 §5) supervises the model with
judgment outcomes while \emph{excluding answer traces}: the model cannot
copy the unfolding; it must induce the computation from definitions and
compile it. The compiled product is intuition --- the reasoning fast
path.

\textbf{Hard to obtain.} Compilation is not free. Acquisition is
config-dependent: 3 epochs for the simple suite, 8 for the full
logic-and-arithmetic suite (exp53). Until compilation completes, the
intuition channel does not generalize (exp53: epochs 1--3 at 0.1--24\%).
This is the \emph{obtain} side of the decoupling.

\textbf{Easy to generalize.} Once compiled, the intuition channel
generalizes by \emph{skipping construction steps}: it jumps from input
to output over the intermediate unfolding. Zero decay out to 2000
digits, radix 60, and 100\% digit anonymization (exp80). This is the
\emph{generalize} side.

\textbf{Decoupling (C4).} The two facts --- hard to obtain, easy to skip
--- are not in tension. Obtaining requires compiling the structure
(training); generalizing uses the compiled structure (inference). The
theory of P1 explains both: definition completeness + relational
sufficiency determine whether compilation succeeds; step-skipping is
what compiled structure does on new inputs.

\subsection{4 Experiments}\label{experiments}

All via the official run\_exp judge path (full-sequence reconstruction).
Model: exp02\_supervised\_s2 (2-layer, dim-64, \textless1 MB, judgment
1.000 on the full suite).

\subsubsection{4.1 Three-channel equivalence
(C2)}\label{three-channel-equivalence-c2}

\begin{longtable}[]{@{}llll@{}}
\toprule\noalign{}
Input & Construction & Formal & Intuition \\
\midrule\noalign{}
\endhead
\bottomrule\noalign{}
\endlastfoot
12+34 & ✓ & ✓ & ✓ \\
123456789+987654321 & ✓ & ✓ & ✓ \\
999\ldots9 (20 digits)+1 & ✓ & ✓ & ✓ \\
\textbf{999\ldots9 (2000 digits)+1} & ✓ & ✓ & \textbf{1.000} \\
\end{longtable}

Construction = formal exactly on full output sequences (2/9/20/2000
digits). Intuition scores 1.000 on the same inputs. The three channels
realize the same computation.

\subsubsection{4.2 Intuition carries structure, not symbols
(C3)}\label{intuition-carries-structure-not-symbols-c3}

\begin{itemize}
\tightlist
\item
  \textbf{Symbol permutation (exp41)}: 26080 digit renamings during
  training, tested on original symbols --- invariance 0.996→0.987. The
  intuition channel realizes structural relations
  (successor/iteration/carry), not digit signifiers.
\item
  \textbf{Step-skipping zero decay (exp80)}: 2000 digits / radix 60 /
  100\% anonymization / joint, all 1.000. The compiled structure
  extrapolates without decay.
\item
  \textbf{Machine form (exp90)}: 20-digit addition compiles 82
  construction tokens into 5 intuition tokens --- the skip is visible in
  the token stream. Honest caveat: the atomic token is outside the model
  vocab (pad id 0); measurement is the sequence-length effect, a lower
  bound.
\end{itemize}

\subsubsection{4.3 Acquisition--generalization decoupling
(C4)}\label{acquisitiongeneralization-decoupling-c4}

\begin{longtable}[]{@{}llllll@{}}
\toprule\noalign{}
Epochs & 1 & 2 & 3 & 5 & 8 \\
\midrule\noalign{}
\endhead
\bottomrule\noalign{}
\endlastfoot
Judgment acc (full suite) & 0.001 & 0.214 & 0.244 & 0.781 & 0.996 \\
\end{longtable}

Acquisition is monotone but config-dependent (8 epochs for the full
suite). The same converged model then extrapolates with zero decay
(§4.2). The ``hard to obtain'' and ``easy to skip'' facts coexist; they
describe obtain (training) and generalize (inference), respectively.

\subsection{5 Discussion}\label{discussion}

\textbf{What the equivalence does not say.} We do not claim the three
philosophical positions are identical in their epistemology; we claim
they are equivalent as \emph{execution channels of one system}. The
historical disagreement concerned the nature of mathematical truth; our
claim concerns the architecture of a computing system that realizes
mathematical objects. The philosophical positions are models of what a
mathematical object is; we show one system can realize all three models
of the same object consistently.

\textbf{Anti-Platonism.} Intuition is compiled construction, not an a
priori faculty. If mathematical structure were innate (Platonism /
nativism), then definitionally-impoverished concepts would still
generalize; they do not (P1 §8: an impoverished vocabulary requires an
order of magnitude more training; genuinely under-defined tokens never
generalize). The intuition channel's competence is earned by
compilation, not endowed.

\textbf{Positional vs constructive content.} Structuralism says an
object is its position; the intuition channel locates the object in the
learned structure. Intuitionism says an object is its construction; the
construction channel unfolds it. The equivalence of the channels shows
the two descriptions coincide --- for objects whose definitions are
complete. This is the machine-level form of the old
structuralist--intuitionist rapprochement.

\subsection{6 Related Work}\label{related-work}

\begin{itemize}
\tightlist
\item
  \textbf{Computational Trinitarianism} (Curry--Howard): programs =
  proofs = types. We extend the three-way identification to execution
  channels of one neural system.
\item
  \textbf{Constructive mathematics / type theory} (Martin-Löf 1984;
  Bishop): existence = constructibility. Our construction channel is its
  machine form.
\item
  \textbf{Structuralism / Bourbaki; Gärdenfors conceptual spaces}:
  object = position in a structure/geometry. Our intuition channel is
  its neural realization.
\item
  \textbf{Formalism / Hilbert; rewriting systems}: object = sign
  manipulation. Our formal channel is the definition-driven engine.
\item
  \textbf{Dual-process cognition} (Kahneman): System 1/2. Our
  intuition/reasoning separation is the System 1/2 split with a
  compilation account of how System 1 is built.
\item
  \textbf{Dual-systems in math cognition} (e.g.~number sense /
  deliberate arithmetic): our three channels refine the two-system
  picture with a verifiable equivalence.
\end{itemize}

\subsection{6.5 Mathematical-domain validation: the Riemann-direction
case}\label{mathematical-domain-validation-the-riemann-direction-case}

The three-channel account is not confined to arithmetic. In a companion
formalization (Lean 4 / mathlib, claims C011--C025, all PROVED, novelty
KNOWN, no sorry), the intuition chain ``the complex axis is a
high-dimensional projection of −1; primes land on translated integer
points; 1/2 is the inversion--translation symmetry center; the complex
axis is curled'' was carried through all three channels:

\begin{itemize}
\tightlist
\item
  \textbf{Construction} (definition-driven): the ComplexAxis algebra (a
  + bJ, J² = −1), projection proj⟨a,b⟩ = a, lifting, basepoint drift;
\item
  \textbf{Formal} (rewriting): Lean 4 proofs of projection
  non-preservation, prime-circle structure (8 lattice points as a single
  orbit under rotation × conjugation, via Gaussian-integer UFD),
  Euler-product convergence (Re(s) \textgreater{} 1), and the
  critical-line geometry;
\item
  \textbf{Intuition} (compiled structure): the same objects are reached
  directly by the intuition chain, skipping textbook derivation.
\end{itemize}

\textbf{Projection recovery theorem} (\texttt{proj\_not\_recoverable} /
\texttt{proj\_recoverable\_symmetry}, Lean-proved): \emph{lost structure
is not recoverable; symmetry directions are recoverable.} Under
projection, the imaginary-axis direction (i vs −i) is lost and cannot be
recovered, while the real-axis sign symmetry (r vs −r) is preserved.
This is the mathematical counterpart of EXP-61h (§4): the intuition
channel is position-insensitive (it drops positional detail like
projection drops the J direction) yet preserves semantic symmetry ---
matching ``scrambled letters do not impede understanding; deliberate
logic on the scramble loses meaning'' (G5b, P1). Projection is the
mechanism of the fast path: dropping irrelevant structure irrecoverably
is what makes the direct route cheap.

\textbf{Token economy of intuition-guided formalization} (methodological
data): the full C011--C025 formalization consumed ≈700k tokens, of which
99.2\% were context-cache re-transmission; net new content \textless{}
1\%. The intuition chain hit known structures directly (heap/torsor,
sums of two squares, circle inversion, functional-equation symmetry),
avoiding textbook derivation --- an efficiency datum for G5 at the token
level, independent of the arithmetic experiments in §4.

\subsection{7 Limitations}\label{limitations}

\begin{enumerate}
\def\labelenumi{\arabic{enumi}.}
\tightlist
\item
  The equivalence is demonstrated on the supported calculus (formalized
  definitions), not on all of mathematics.
\item
  The intuition channel is statistical, not guaranteed; equivalence is
  verified (held-out) and falls back to construction, not proven.
\item
  The atomization result (exp90) is a sequence-length lower bound, not a
  claim of learned atomic semantics.
\end{enumerate}

\subsection{8 Conclusion}\label{conclusion}

The three positions of the foundations dispute --- construction,
structure, sign --- are realized as three equivalent execution channels
of one neural system. Construction = formal exactly; intuition agrees at
1.000; the intuition channel carries structure, not symbols; and it is
hard to obtain but easy to generalize by skipping steps. Intuition is
compiled construction. A century-old split in the foundations of
mathematics is, at least for a formalized calculus, an architectural
fact about how many ways one system can realize the same object --- and
all of them agree.

\subsection{Appendix}\label{appendix}

\begin{itemize}
\tightlist
\item
  A. Full experiment tables: docs/paper\_data/results.csv.
\item
  B. Relation to P1 propositions (G0, G5, R) and P2 compilation.
\item
  C. Formal-verification direction (Lean): construction ≡ intuition ≡
  formal, for the supported calculus.
\end{itemize}

\begin{center}\rule{0.5\linewidth}{0.5pt}\end{center}

\subsection{Unified Conclusion (added
2026-08-11)}\label{unified-conclusion-added-2026-08-11}

Once, under a unified formal attention-syntactic structure, a neural
network ultimately, inevitably, and necessarily clearly, intuitively,
and interpretably observes the boundary between intuition and
construction, we can with full confidence use the unified
attention-formal language to direct logical construction, use precise
synthetic-CoT logical construction to earn \emph{flashing intuition},
and use distilled intuitive context to exhaustively search attention
forms. Form, construction, and intuition co-adapt within the unified
framework, forming a feedback iteration pipeline of teaching, training,
and reasoning in which they are no longer distinct --- perhaps one of
the faster shortcuts among the many paths toward ASI.

\subsection{Cross-References within the Unified
Framework}\label{cross-references-within-the-unified-framework}

This paper is one leg of the \emph{Unified Framework of Representation,
Logic, and Intuition} (manifesto:
\texttt{Unified\_Framework\_Representation\_Logic\_Intuition.md}). The
three legs co-adapt:

\begin{longtable}[]{@{}
  >{\raggedright\arraybackslash}p{(\linewidth - 4\tabcolsep) * \real{0.3333}}
  >{\raggedright\arraybackslash}p{(\linewidth - 4\tabcolsep) * \real{0.3333}}
  >{\raggedright\arraybackslash}p{(\linewidth - 4\tabcolsep) * \real{0.3333}}@{}}
\toprule\noalign{}
\begin{minipage}[b]{\linewidth}\raggedright
Leg
\end{minipage} & \begin{minipage}[b]{\linewidth}\raggedright
Paper
\end{minipage} & \begin{minipage}[b]{\linewidth}\raggedright
Claim
\end{minipage} \\
\midrule\noalign{}
\endhead
\bottomrule\noalign{}
\endlastfoot
form → construction & P1 (this paper) & correct syntax → extreme
training speed; correct construction → extreme generalization \\
construction → intuition & P0 (three-channel equivalence) & construction
= formal = intuition; intuition is compiled construction \\
construction → intuition (method) & P2 (neural macro compilation) &
compile the slow construction path into the fast reasoning path \\
\end{longtable}

Cyclic closure: form directs construction (P1), construction earns
intuition (P0/P2), intuition searches form (the fast path extrapolates
and re-expresses in the unified form). The unified conclusion --- form
directs logical construction, construction earns flashing intuition,
intuition exhaustively searches forms --- is stated in the manifesto and
in each paper's Unified Conclusion.

\end{document}
